\begin{register}{H}{GPIO\_BT\_SELECT\_REG}{0x{}0000}\label{regdesc:GPIOBTSELECTREG}
\regfield{GPIO\_BT\_SEL}{32}{0}{{0x{}000000}}%
\reglabel{Reset}\regnewline%
\begin{regdesc}\begin{reglist}
\label{fielddesc:GPIOBTSEL}\item [GPIO\_BT\_SEL] Reserved (R/W)
\end{reglist}\end{regdesc}
\end{register}


\begin{register}{H}{GPIO\_OUT\_REG}{0x{}0004}\label{regdesc:GPIOOUTREG}
\regfield{GPIO\_OUT\_DATA\_ORIG}{32}{0}{{0x{}000000}}%
\reglabel{Reset}\regnewline%
\begin{regdesc}\begin{reglist}
\label{fielddesc:GPIOOUTDATAORIG}\item [GPIO\_OUT\_DATA\_ORIG] GPIO0 \verb+~+ 21 and GPIO26 \verb+~+ 31 output values in simple GPIO output mode. The values of bit0 \verb+~+ bit21 correspond to the output values of GPIO0 \verb+~+ 21, and bit26 \verb+~+ bit31 to GPIO26 \verb+~+ 31. Bit22 \verb+~+ bit25 are invalid. (R/W)
\end{reglist}\end{regdesc}
\end{register}


\begin{register}{H}{GPIO\_OUT\_W1TS\_REG}{0x{}0008}\label{regdesc:GPIOOUTW1TSREG}
\regfield{GPIO\_OUT\_W1TS}{32}{0}{{0x{}000000}}%
\reglabel{Reset}\regnewline%
\begin{regdesc}\begin{reglist}
\label{fielddesc:GPIOOUTW1TS}\item [GPIO\_OUT\_W1TS] GPIO0 \verb+~+ 31 output set register. If the value 1 is written to a bit here, the corre- sponding bit in \hyperref[regdesc:GPIOOUTREG]{GPIO\_OUT\_REG} will be set to 1. Recommended operation: use this register to set \hyperref[regdesc:GPIOOUTREG]{GPIO\_OUT\_REG}. (WO)
\end{reglist}\end{regdesc}
\end{register}


\begin{register}{H}{GPIO\_OUT\_W1TC\_REG}{0x{}000C}\label{regdesc:GPIOOUTW1TCREG}
\regfield{GPIO\_OUT\_W1TC}{32}{0}{{0x{}000000}}%
\reglabel{Reset}\regnewline%
\begin{regdesc}\begin{reglist}
\label{fielddesc:GPIOOUTW1TC}\item [GPIO\_OUT\_W1TC] GPIO0 \verb+~+ 31 output clear register. If the value 1 is written to a bit here, the cor- responding bit in \hyperref[regdesc:GPIOOUTREG]{GPIO\_OUT\_REG} will be cleared. Recommended operation: use this register to clear \hyperref[regdesc:GPIOOUTREG]{GPIO\_OUT\_REG}. (WO)
\end{reglist}\end{regdesc}
\end{register}


\begin{register}{H}{GPIO\_OUT1\_REG}{0x{}0010}\label{regdesc:GPIOOUT1REG}
\regfield{(reserved)}{10}{22}{0000000000}%
\regfield{GPIO\_OUT1\_DATA\_ORIG}{22}{0}{{0x{}0000}}%
\reglabel{Reset}\regnewline%
\begin{regdesc}\begin{reglist}
\label{fielddesc:GPIOOUT1DATAORIG}\item [GPIO\_OUT1\_DATA\_ORIG] GPIO32 \verb+~+ 48 output value in simple GPIO output mode. The values of bit0 \verb+~+ bit16 correspond to GPIO32 \verb+~+ GPIO48. Bit17 \verb+~+ bit21 are invalid. (R/W)
\end{reglist}\end{regdesc}
\end{register}


\begin{register}{H}{GPIO\_OUT1\_W1TS\_REG}{0x{}0014}\label{regdesc:GPIOOUT1W1TSREG}
\regfield{(reserved)}{10}{22}{0000000000}%
\regfield{GPIO\_OUT1\_W1TS}{22}{0}{{0x{}0000}}%
\reglabel{Reset}\regnewline%
\begin{regdesc}\begin{reglist}
\label{fielddesc:GPIOOUT1W1TS}\item [GPIO\_OUT1\_W1TS] GPIO32 \verb+~+ 48 output value set register. If the value 1 is written to a bit here, the corresponding bit in \hyperref[regdesc:GPIOOUT1REG]{GPIO\_OUT1\_REG} will be set to 1. Recommended operation: use this register to set \hyperref[regdesc:GPIOOUT1REG]{GPIO\_OUT1\_REG}. (WO)
\end{reglist}\end{regdesc}
\end{register}


\begin{register}{H}{GPIO\_OUT1\_W1TC\_REG}{0x{}0018}\label{regdesc:GPIOOUT1W1TCREG}
\regfield{(reserved)}{10}{22}{0000000000}%
\regfield{GPIO\_OUT1\_W1TC}{22}{0}{{0x{}0000}}%
\reglabel{Reset}\regnewline%
\begin{regdesc}\begin{reglist}
\label{fielddesc:GPIOOUT1W1TC}\item [GPIO\_OUT1\_W1TC] GPIO32 \verb+~+ 48 output value clear register. If the value 1 is written to a bit here, the corresponding bit in \hyperref[regdesc:GPIOOUT1REG]{GPIO\_OUT1\_REG} will be cleared. Recommended operation: use this register to clear \hyperref[regdesc:GPIOOUT1REG]{GPIO\_OUT1\_REG}. (WO)
\end{reglist}\end{regdesc}
\end{register}


\begin{register}{H}{GPIO\_SDIO\_SELECT\_REG}{0x{}001C}\label{regdesc:GPIOSDIOSELECTREG}
\regfield{(reserved)}{24}{8}{000000000000000000000000}%
\regfield{GPIO\_SDIO\_SEL}{8}{0}{{0x{}0}}%
\reglabel{Reset}\regnewline%
\begin{regdesc}\begin{reglist}
\label{fielddesc:GPIOSDIOSEL}\item [GPIO\_SDIO\_SEL] Reserved (R/W)
\end{reglist}\end{regdesc}
\end{register}


\begin{register}{H}{GPIO\_ENABLE\_REG}{0x{}0020}\label{regdesc:GPIOENABLEREG}
\regfield{GPIO\_ENABLE\_DATA}{32}{0}{{0x{}000000}}%
\reglabel{Reset}\regnewline%
\begin{regdesc}\begin{reglist}
\label{fielddesc:GPIOENABLEDATA}\item [GPIO\_ENABLE\_DATA] GPIO0\verb+~+31 output enable register. (R/W)
\end{reglist}\end{regdesc}
\end{register}


\begin{register}{H}{GPIO\_ENABLE\_W1TS\_REG}{0x{}0024}\label{regdesc:GPIOENABLEW1TSREG}
\regfield{GPIO\_ENABLE\_W1TS}{32}{0}{{0x{}000000}}%
\reglabel{Reset}\regnewline%
\begin{regdesc}\begin{reglist}
\label{fielddesc:GPIOENABLEW1TS}\item [GPIO\_ENABLE\_W1TS] GPIO0 \verb+~+ 31 output enable set register. If the value 1 is written to a bit here, the corresponding bit in \hyperref[regdesc:GPIOENABLEREG]{GPIO\_ENABLE\_REG} will be set to 1. Recommended operation: use this register to set \hyperref[regdesc:GPIOENABLEREG]{GPIO\_ENABLE\_REG}. (WO)
\end{reglist}\end{regdesc}
\end{register}


\begin{register}{H}{GPIO\_ENABLE\_W1TC\_REG}{0x{}0028}\label{regdesc:GPIOENABLEW1TCREG}
\regfield{GPIO\_ENABLE\_W1TC}{32}{0}{{0x{}000000}}%
\reglabel{Reset}\regnewline%
\begin{regdesc}\begin{reglist}
\label{fielddesc:GPIOENABLEW1TC}\item [GPIO\_ENABLE\_W1TC] GPIO0 \verb+~+ 31 output enable clear register. If the value 1 is written to a bit here, the corresponding bit in \hyperref[regdesc:GPIOENABLEREG]{GPIO\_ENABLE\_REG} will be cleared. Recommended operation: use this register to clear \hyperref[regdesc:GPIOENABLEREG]{GPIO\_ENABLE\_REG}. (WO)
\end{reglist}\end{regdesc}
\end{register}


\begin{register}{H}{GPIO\_ENABLE1\_REG}{0x{}002C}\label{regdesc:GPIOENABLE1REG}
\regfield{(reserved)}{10}{22}{0000000000}%
\regfield{GPIO\_ENABLE1\_DATA}{22}{0}{{0x{}0000}}%
\reglabel{Reset}\regnewline%
\begin{regdesc}\begin{reglist}
\label{fielddesc:GPIOENABLE1DATA}\item [GPIO\_ENABLE1\_DATA] GPIO32 \verb+~+ 48 output enable register. (R/W)
\end{reglist}\end{regdesc}
\end{register}


\begin{register}{H}{GPIO\_ENABLE1\_W1TS\_REG}{0x{}0030}\label{regdesc:GPIOENABLE1W1TSREG}
\regfield{(reserved)}{10}{22}{0000000000}%
\regfield{GPIO\_ENABLE1\_W1TS}{22}{0}{{0x{}0000}}%
\reglabel{Reset}\regnewline%
\begin{regdesc}\begin{reglist}
\label{fielddesc:GPIOENABLE1W1TS}\item [GPIO\_ENABLE1\_W1TS] GPIO32 \verb+~+ 48 output enable set register. If the value 1 is written to a bit here, the corresponding bit in \hyperref[regdesc:GPIOENABLE1REG]{GPIO\_ENABLE1\_REG} will be set to 1. Recommended operation: use this register to set \hyperref[regdesc:GPIOENABLE1REG]{GPIO\_ENABLE1\_REG}. (WO)
\end{reglist}\end{regdesc}
\end{register}


\begin{register}{H}{GPIO\_ENABLE1\_W1TC\_REG}{0x{}0034}\label{regdesc:GPIOENABLE1W1TCREG}
\regfield{(reserved)}{10}{22}{0000000000}%
\regfield{GPIO\_ENABLE1\_W1TC}{22}{0}{{0x{}0000}}%
\reglabel{Reset}\regnewline%
\begin{regdesc}\begin{reglist}
\label{fielddesc:GPIOENABLE1W1TC}\item [GPIO\_ENABLE1\_W1TC] GPIO32 \verb+~+ 48 output enable clear register. If the value 1 is written to a bit here, the corresponding bit in \hyperref[regdesc:GPIOENABLE1REG]{GPIO\_ENABLE1\_REG} will be cleared. Recommended operation: use this register to clear \hyperref[regdesc:GPIOENABLE1REG]{GPIO\_ENABLE1\_REG}. (WO)
\end{reglist}\end{regdesc}
\end{register}


\begin{register}{H}{GPIO\_STRAP\_REG}{0x{}0038}\label{regdesc:GPIOSTRAPREG}
\regfield{(reserved)}{26}{6}{00000000000000000000000000}%
\regfield{GPIO\_STRAPPING}{6}{0}{{0}}%
\reglabel{Reset}\regnewline%
\begin{regdesc}\begin{reglist}
\label{fielddesc:GPIOSTRAPPING}\item [GPIO\_STRAPPING] Represents the values of corresponding GPIO Strapping pins. \\
bit0 \verb+~+ bit1: Reserved\\
bit2: GPIO46\\
bit3: GPIO0\\
bit4: GPIO45\\
bit5: GPIO3\\(RO)
\end{reglist}\end{regdesc}
\end{register}


\begin{register}{H}{GPIO\_IN\_REG}{0x{}003C}\label{regdesc:GPIOINREG}
\regfield{GPIO\_IN\_DATA\_NEXT}{32}{0}{{0}}%
\reglabel{Reset}\regnewline%
\begin{regdesc}\begin{reglist}
\label{fielddesc:GPIOINDATANEXT}\item [GPIO\_IN\_DATA\_NEXT] GPIO0 \verb+~+ 31 input value. Each bit represents a pin input value, 1 for high level and 0 for low level. (RO)
\end{reglist}\end{regdesc}
\end{register}


\begin{register}{H}{GPIO\_IN1\_REG}{0x{}0040}\label{regdesc:GPIOIN1REG}
\regfield{(reserved)}{10}{22}{0000000000}%
\regfield{GPIO\_IN\_DATA1\_NEXT}{22}{0}{{0}}%
\reglabel{Reset}\regnewline%
\begin{regdesc}\begin{reglist}
\label{fielddesc:GPIOINDATA1NEXT}\item [GPIO\_IN\_DATA1\_NEXT] GPIO32 \verb+~+ 48 input value. Each bit represents a pin input value. (RO)
\end{reglist}\end{regdesc}
\end{register}


\begin{register}{H}{GPIO\_PIN\regindex{n}\_REG (\regindex{n}: 0-48)}{0x{}0074+0x{}4*\regindex{n}}\label{regdesc:GPIOPINNREG}
\regfield{(reserved)}{14}{18}{00000000000000}%
\regfield{GPIO\_PIN\regindex{n}\_INT\_ENA}{5}{13}{{0x{}0}}%
\regfield{GPIO\_PIN\regindex{n}\_CONFIG}{2}{11}{{0x{}0}}%
\regfield{GPIO\_PIN\regindex{n}\_WAKEUP\_ENABLE}{1}{10}{{0}}%
\regfield{GPIO\_PIN\regindex{n}\_INT\_TYPE}{3}{7}{{0x{}0}}%
\regfield{(reserved)}{2}{5}{00}%
\regfield{GPIO\_PIN\regindex{n}\_SYNC1\_BYPASS}{2}{3}{{0x{}0}}%
\regfield{GPIO\_PIN\regindex{n}\_PAD\_DRIVER}{1}{2}{{0}}%
\regfield{GPIO\_PIN\regindex{n}\_SYNC2\_BYPASS}{2}{0}{{0x{}0}}%
\reglabel{Reset}\regnewline%
\begin{regdesc}\begin{reglist}
\label{fielddesc:GPIOPINNSYNC2BYPASS}\item [GPIO\_PIN\regindex{n}\_SYNC2\_BYPASS] For the second stage synchronization, GPIO input data can be syn- chronized on either edge of the APB clock. 0: no synchronization; 1: synchronized on falling edge; 2 and 3: synchronized on rising edge. (R/W)
\label{fielddesc:GPIOPINNPADDRIVER}\item [GPIO\_PIN\regindex{n}\_PAD\_DRIVER] Pin driver selection. 0: normal output; 1: open drain output. (R/W)
\label{fielddesc:GPIOPINNSYNC1BYPASS}\item [GPIO\_PIN\regindex{n}\_SYNC1\_BYPASS] For the first stage synchronization, GPIO input data can be synchro- nized on either edge of the APB clock. 0: no synchronization; 1: synchronized on falling edge; 2 and 3: synchronized on rising edge. (R/W)
\label{fielddesc:GPIOPINNINTTYPE}\item [GPIO\_PIN\regindex{n}\_INT\_TYPE] Interrupt type selection. 0: GPIO interrupt disabled; 1: rising edge trigger; 2: falling edge trigger; 3: any edge trigger; 4: low level trigger; 5: high level trigger. (R/W) 
\label{fielddesc:GPIOPINNWAKEUPENABLE}\item [GPIO\_PIN\regindex{n}\_WAKEUP\_ENABLE] GPIO wake-up enable bit, only wakes up the CPU from Light-sleep. (R/W)
\label{fielddesc:GPIOPINNCONFIG}\item [GPIO\_PIN\regindex{n}\_CONFIG] Reserved (R/W)
\label{fielddesc:GPIOPINNINTENA}\item [GPIO\_PIN\regindex{n}\_INT\_ENA] Interrupt enable bits. bit13: CPU interrupt enabled; bit14: CPU non-maskable interrupt enabled. (R/W)
\end{reglist}\end{regdesc}
\end{register}


\begin{register}{H}{GPIO\_FUNC\regindex{y}\_IN\_SEL\_CFG\_REG (\regindex{y}: 0-255)}{0x{}0154+0x{}4*\regindex{y}}\label{regdesc:GPIOFUNCNINSELCFGREG}
\regfield{(reserved)}{24}{8}{000000000000000000000000}%
\regfield{GPIO\_SIG\regindex{y}\_IN\_SEL}{1}{7}{{0}}%
\regfield{GPIO\_FUNC\regindex{y}\_IN\_INV\_SEL}{1}{6}{{0}}%
\regfield{GPIO\_FUNC\regindex{y}\_IN\_SEL}{6}{0}{{0x{}0}}%
\reglabel{Reset}\regnewline%
\begin{regdesc}\begin{reglist}
\label{fielddesc:GPIOFUNCNINSEL}\item [GPIO\_FUNC\regindex{y}\_IN\_SEL] Selection control for peripheral input signal \regindex{Y}, selects a pin from the 48 GPIO matrix pins to connect this input signal. Or selects 0x{}38 for a constantly high input or 0x{}3C for a constantly low input. (R/W)
\label{fielddesc:GPIOFUNCNININVSEL}\item [GPIO\_FUNC\regindex{y}\_IN\_INV\_SEL] 1: Invert the input value; 0: Do not invert the input value. (R/W)
\label{fielddesc:GPIOSIGNINSEL}\item [GPIO\_SIG\regindex{y}\_IN\_SEL] Bypass GPIO matrix. 1: route signals via GPIO matrix, 0: connect signals directly to peripheral configured in IO MUX. (R/W)
\end{reglist}\end{regdesc}
\end{register}


\begin{register}{H}{GPIO\_FUNC\regindex{x}\_OUT\_SEL\_CFG\_REG (\regindex{x}: 0-48)}{0x{}0554+0x{}4*\regindex{x}}\label{regdesc:GPIOFUNCNOUTSELCFGREG}
\regfield{(reserved)}{20}{12}{00000000000000000000}%
\regfield{GPIO\_FUNC\regindex{x}\_OEN\_INV\_SEL}{1}{11}{{0}}%
\regfield{GPIO\_FUNC\regindex{x}\_OEN\_SEL}{1}{10}{{0}}%
\regfield{GPIO\_FUNC\regindex{x}\_OUT\_INV\_SEL}{1}{9}{{0}}%
\regfield{GPIO\_FUNC\regindex{x}\_OUT\_SEL}{9}{0}{{0x{}100}}%
\reglabel{Reset}\regnewline%
\begin{regdesc}\begin{reglist}
\label{fielddesc:GPIOFUNCNOUTSEL}\item [GPIO\_FUNC\regindex{x}\_OUT\_SEL] Selection control for GPIO output \regindex{X}. If a value \regindex{Y} (0<=Y<256) is written to this field, the peripheral output signal \regindex{Y} will be connected to GPIO output \regindex{X}. If a value 256 is written to this field, bit \regindex{X} of \hyperref[regdesc:GPIOOUTREG]{GPIO\_OUT\_REG}/\hyperref[regdesc:GPIOOUT1REG]{GPIO\_OUT1\_REG} and \hyperref[regdesc:GPIOENABLEREG]{GPIO\_ENABLE\_REG}/\hyperref[regdesc:GPIOENABLE1REG]{GPIO\_ENABLE1\_REG} will be selected as the output value and output enable. (R/W)
\label{fielddesc:GPIOFUNCNOUTINVSEL}\item [GPIO\_FUNC\regindex{x}\_OUT\_INV\_SEL] 0: Do not invert the output value; 1: Invert the output value. (R/W)
\label{fielddesc:GPIOFUNCNOENSEL}\item [GPIO\_FUNC\regindex{x}\_OEN\_SEL] 0: Use output enable signal from peripheral; 1: Force the output enable signal to be sourced from \hyperref[regdesc:GPIOENABLEREG]{GPIO\_ENABLE\_REG}[\regindex{x}]. (R/W)
\label{fielddesc:GPIOFUNCNOENINVSEL}\item [GPIO\_FUNC\regindex{n}\_OEN\_INV\_SEL] 0: Do not invert the output enable signal; 1: Invert the output enable signal. (R/W)
\end{reglist}\end{regdesc}
\end{register}


\begin{register}{H}{GPIO\_CLOCK\_GATE\_REG}{0x{}062C}\label{regdesc:GPIOCLOCKGATEREG}
\regfield{(reserved)}{31}{1}{0000000000000000000000000000000}%
\regfield{GPIO\_CLK\_EN}{1}{0}{{1}}%
\reglabel{Reset}\regnewline%
\begin{regdesc}\begin{reglist}
\label{fielddesc:GPIOCLKEN}\item [GPIO\_CLK\_EN] Clock gating enable bit. If set to 1, the clock is free running. (R/W)
\end{reglist}\end{regdesc}
\end{register}


\begin{register}{H}{GPIO\_STATUS\_REG}{0x{}0044}\label{regdesc:GPIOSTATUSREG}
\regfield{GPIO\_STATUS\_INTERRUPT}{32}{0}{{0x{}000000}}%
\reglabel{Reset}\regnewline%
\begin{regdesc}\begin{reglist}
\label{fielddesc:GPIOSTATUSINTERRUPT}\item [GPIO\_STATUS\_INTERRUPT] GPIO0 \verb+~+ 31 interrupt status register. (R/W)
\end{reglist}\end{regdesc}
\end{register}


\begin{register}{H}{GPIO\_STATUS1\_REG}{0x{}0050}\label{regdesc:GPIOSTATUS1REG}
\regfield{(reserved)}{10}{22}{0000000000}%
\regfield{GPIO\_STATUS1\_INTERRUPT}{22}{0}{{0x{}0000}}%
\reglabel{Reset}\regnewline%
\begin{regdesc}\begin{reglist}
\label{fielddesc:GPIOSTATUS1INTERRUPT}\item [GPIO\_STATUS1\_INTERRUPT] GPIO32 \verb+~+ 48 interrupt status register. (R/W)
\end{reglist}\end{regdesc}
\end{register}


\begin{register}{H}{GPIO\_CPU\_INT\_REG}{0x{}005C}\label{regdesc:GPIOPCPUINTREG}
\regfield{GPIO\_CPU\_INT}{32}{0}{{0x{}000000}}%
\reglabel{Reset}\regnewline%
\begin{regdesc}\begin{reglist}
\label{fielddesc:GPIOPROCPUINT}\item [GPIO\_CPU\_INT] GPIO0 \verb+~+ 31 CPU interrupt status. This interrupt status is corresponding to the bit in \hyperref[regdesc:GPIOSTATUSREG]{GPIO\_STATUS\_REG} when assert (high) enable signal (bit13 of \hyperref[regdesc:GPIOPINNREG]{GPIO\_PIN\regindex{n}\_REG}). (RO)
\end{reglist}\end{regdesc}
\end{register}


\begin{register}{H}{GPIO\_CPU\_NMI\_INT\_REG}{0x{}0060}\label{regdesc:GPIOPCPUNMIINTREG}
\regfield{GPIO\_CPU\_NMI\_INT}{32}{0}{{0x{}000000}}%
\reglabel{Reset}\regnewline%
\begin{regdesc}\begin{reglist}
\label{fielddesc:GPIOPROCPUNMIINT}\item [GPIO\_CPU\_NMI\_INT] GPIO0 \verb+~+ 31 CPU non-maskable interrupt status. This interrupt status is corresponding to the bit in \hyperref[regdesc:GPIOSTATUSREG]{GPIO\_STATUS\_REG} when assert (high) enable signal (bit 14 of \hyperref[regdesc:GPIOPINNREG]{GPIO\_PIN\regindex{n}\_REG}). (RO)
\end{reglist}\end{regdesc}
\end{register}


\begin{register}{H}{GPIO\_CPU\_INT1\_REG}{0x{}0068}\label{regdesc:GPIOPCPUINT1REG}
\regfield{(reserved)}{10}{22}{0000000000}%
\regfield{GPIO\_CPU1\_INT}{22}{0}{{0x{}0000}}%
\reglabel{Reset}\regnewline%
\begin{regdesc}\begin{reglist}
\label{fielddesc:GPIOPROCPU1INT}\item [GPIO\_CPU1\_INT] GPIO32 \verb+~+ 48 CPU interrupt status. This interrupt status is corresponding to the bit in \hyperref[regdesc:GPIOSTATUS1REG]{GPIO\_STATUS1\_REG} when assert (high) enable signal (bit 13 of \hyperref[regdesc:GPIOPINNREG]{GPIO\_PIN\regindex{n}\_REG}). (RO)
\end{reglist}\end{regdesc}
\end{register}


\begin{register}{H}{GPIO\_CPU\_NMI\_INT1\_REG}{0x{}006C}\label{regdesc:GPIOPCPUNMIINT1REG}
\regfield{(reserved)}{10}{22}{0000000000}%
\regfield{GPIO\_CPU\_NMI1\_INT}{22}{0}{{0x{}0000}}%
\reglabel{Reset}\regnewline%
\begin{regdesc}\begin{reglist}
\label{fielddesc:GPIOPROCPUNMI1INT}\item [GPIO\_CPU\_NMI1\_INT] GPIO32 \verb+~+ 48 CPU non-maskable interrupt status. This interrupt status is corresponding to bit in \hyperref[regdesc:GPIOSTATUS1REG]{GPIO\_STATUS1\_REG} when assert (high) enable signal (bit 14 of \hyperref[regdesc:GPIOPINNREG]{GPIO\_PIN\regindex{n}\_REG}). (RO)
\end{reglist}\end{regdesc}
\end{register}


\begin{register}{H}{GPIO\_STATUS\_W1TS\_REG}{0x{}0048}\label{regdesc:GPIOSTATUSW1TSREG}
\regfield{GPIO\_STATUS\_W1TS}{32}{0}{{0x{}000000}}%
\reglabel{Reset}\regnewline%
\begin{regdesc}\begin{reglist}
\label{fielddesc:GPIOSTATUSW1TS}\item [GPIO\_STATUS\_W1TS] GPIO0 \verb+~+ 31 interrupt status set register. If the value 1 is written to a bit here, the corresponding bit in \hyperref[fielddesc:GPIOSTATUSINTERRUPT]{GPIO\_STATUS\_INTERRUPT} will be set to 1. Recommended operation: use this register to set \hyperref[fielddesc:GPIOSTATUSINTERRUPT]{GPIO\_STATUS\_INTERRUPT}. (WO)
\end{reglist}\end{regdesc}
\end{register}


\begin{register}{H}{GPIO\_STATUS\_W1TC\_REG}{0x{}004C}\label{regdesc:GPIOSTATUSW1TCREG}
\regfield{GPIO\_STATUS\_W1TC}{32}{0}{{0x{}000000}}%
\reglabel{Reset}\regnewline%
\begin{regdesc}\begin{reglist}
\label{fielddesc:GPIOSTATUSW1TC}\item [GPIO\_STATUS\_W1TC] GPIO0 \verb+~+ 31 interrupt status clear register. If the value 1 is written to a bit here, the corresponding bit in \hyperref[fielddesc:GPIOSTATUSINTERRUPT]{GPIO\_STATUS\_INTERRUPT} will be cleared. Recommended operation: use this register to clear \hyperref[fielddesc:GPIOSTATUSINTERRUPT]{GPIO\_STATUS\_INTERRUPT}. (WO)
\end{reglist}\end{regdesc}
\end{register}


\begin{register}{H}{GPIO\_STATUS1\_W1TS\_REG}{0x{}0054}\label{regdesc:GPIOSTATUS1W1TSREG}
\regfield{(reserved)}{10}{22}{0000000000}%
\regfield{GPIO\_STATUS1\_W1TS}{22}{0}{{0x{}0000}}%
\reglabel{Reset}\regnewline%
\begin{regdesc}\begin{reglist}
\label{fielddesc:GPIOSTATUS1W1TS}\item [GPIO\_STATUS1\_W1TS] GPIO32 \verb+~+ 48 interrupt status set register. If the value 1 is written to a bit here, the corresponding bit in \hyperref[regdesc:GPIOSTATUS1REG]{GPIO\_STATUS1\_REG} will be set to 1. Recommended operation: use this register to set \hyperref[regdesc:GPIOSTATUS1REG]{GPIO\_STATUS1\_REG}. (WO)
\end{reglist}\end{regdesc}
\end{register}


\begin{register}{H}{GPIO\_STATUS1\_W1TC\_REG}{0x{}0058}\label{regdesc:GPIOSTATUS1W1TCREG}
\regfield{(reserved)}{10}{22}{0000000000}%
\regfield{GPIO\_STATUS1\_W1TC}{22}{0}{{0x{}0000}}%
\reglabel{Reset}\regnewline%
\begin{regdesc}\begin{reglist}
\label{fielddesc:GPIOSTATUS1W1TC}\item [GPIO\_STATUS1\_W1TC] GPIO32 \verb+~+ 48 interrupt status clear register. If the value 1 is written to a bit here, the corresponding bit in \hyperref[regdesc:GPIOSTATUS1REG]{GPIO\_STATUS1\_REG} will be cleared. Recommended operation: use this register to clear \hyperref[regdesc:GPIOSTATUS1REG]{GPIO\_STATUS1\_REG}. (WO)
\end{reglist}\end{regdesc}
\end{register}


\begin{register}{H}{GPIO\_STATUS\_NEXT\_REG}{0x{}014C}\label{regdesc:GPIOSTATUSNEXTREG}
\regfield{GPIO\_STATUS\_INTERRUPT\_NEXT}{32}{0}{{0x{}000000}}%
\reglabel{Reset}\regnewline%
\begin{regdesc}\begin{reglist}
\label{fielddesc:GPIOSTATUSINTERRUPTNEXT}\item [GPIO\_STATUS\_INTERRUPT\_NEXT] Interrupt source signal of GPIO0 \verb+~+ 31, could be rising edge interrupt, falling edge interrupt, level sensitive interrupt and any edge interrupt. (RO)
\end{reglist}\end{regdesc}
\end{register}


\begin{register}{H}{GPIO\_STATUS\_NEXT1\_REG}{0x{}0150}\label{regdesc:GPIOSTATUSNEXT1REG}
\regfield{(reserved)}{10}{22}{0000000000}%
\regfield{GPIO\_STATUS1\_INTERRUPT\_NEXT}{22}{0}{{0x{}0000}}%
\reglabel{Reset}\regnewline%
\begin{regdesc}\begin{reglist}
\label{fielddesc:GPIOSTATUS1INTERRUPTNEXT}\item [GPIO\_STATUS1\_INTERRUPT\_NEXT] Interrupt source signal of GPIO32 \verb+~+ 48. (RO)
\end{reglist}\end{regdesc}
\end{register}


\begin{register}{H}{GPIO\_DATE\_REG}{0x{}06FC}\label{regdesc:GPIOREGDATEREG}
\regfield{(reserved)}{4}{28}{0000}%
\regfield{GPIO\_DATE}{28}{0}{{0x{}1907040}}%
\reglabel{Reset}\regnewline%
\begin{regdesc}\begin{reglist}
\label{fielddesc:GPIODATE}\item [GPIO\_DATE] Version control register (R/W)
\end{reglist}\end{regdesc}
\end{register}


